\section{Background and Motivation: Learned Cache Eviction}
\label{sec:background}

At each full-cache miss, a caching system must evict one resident object before admitting the requested one. Classical policies such as LRU operationalize this choice through simple recency structure, while learned or learning-augmented approaches attempt to predict which eviction will lead to lower future miss cost. In practice, these approaches differ in feature sets, prediction targets, and evaluation methodology.

The main difficulty for a benchmark is that victim selection is structured, not i.i.d. A model does not act on one object in isolation; it acts within a decision set whose size, feature spread, and ambiguity vary across traces and capacities. A useful benchmark therefore preserves decision boundaries instead of flattening them away.

This observation motivates the core \lafc\ design choice: represent supervision at the candidate-within-decision level. Each row corresponds to one candidate victim at one full-cache eviction decision, and the label records the finite-horizon counterfactual miss cost of forcing that eviction under a fixed continuation policy. This representation allows both pointwise learning and within-decision evaluation, while keeping the supervised object aligned with the actual cache-eviction choice.

From that candidate-level foundation, the benchmark derives two secondary evaluation surfaces. Grouping all rows with the same \texttt{decision\_id} produces a decision view that records candidate count, optimal-set membership, tie structure, and regret summaries. Sampling candidate pairs from the same decision produces the shipped pairwise view, which turns within-decision comparisons into a simpler classification problem without discarding the original decision boundary.

A single full-cache miss can therefore appear in three complementary forms. If one decision exposes $k$ admissible victims, the release records $k$ candidate rows, one decision-group summary, and up to $k(k-1)/2$ within-decision pairwise comparisons before any sampling cap is applied. The manuscript keeps that task-view progression textual rather than introducing a second overview figure beside the main pipeline diagram.
